\documentclass[11pt]{article}
\usepackage{amsmath,amssymb,amsthm}
\usepackage[margin=1in]{geometry}

\newtheorem{theorem}{Theorem}
\newtheorem{lemma}{Lemma}
\newtheorem{definition}{Definition}

\title{Fixed-\(r\) Route-2 Certificate Target}
\date{2026-05-27}

\begin{document}
\maketitle

\section{Purpose}

This target isolates a certificate-shaped proof obligation from the
Erdos Problem 993 project.  It does not ask for a proof of tree independence
polynomial unimodality.  It asks for formal bridge lemmas that justify the
fixed-\(r\) Route-2 spider-lane certificate.

\section{Notation}

Let \(P_j(x)\) be the independence polynomial of the path on \(j\) vertices.
For fixed \(r \ge 2\), define the spider
\[
  T_{a,r}=S(2^a,r).
\]
Its independence polynomial has the standard spider decomposition
\[
  I(T_{a,r};x)=F_{a,r}(x)+G_{a,r}(x),
  \qquad
  F_{a,r}(x)=P_r(x)(1+2x)^a,\quad
  G_{a,r}(x)=xP_{r-1}(x)(1+x)^a.
\]
Removing one length-\(2\) arm gives
\[
  B_{a,r}=S(2^{a-1},r)
\]
with hub-off part
\[
  F^-_{a,r}(x)=P_r(x)(1+2x)^{a-1}.
\]

For \(H(x)=\sum_k h_kx^k\), write
\[
  \mu_H(\lambda)=\frac{\lambda H'(\lambda)}{H(\lambda)}.
\]
Let \(m=m(a,r)\) be the certified mode of \(I(T_{a,r};x)\), and let
\[
  \lambda=\frac{[x^{m-1}]I(T_{a,r};x)}{[x^m]I(T_{a,r};x)}.
\]
The Route-2 target is
\[
  \mu_{B_{a,r}}(\lambda)\ge m-\frac32.
\]

\section{Certificate Bridge}

For each residue class \(q\in\{0,1,2\}\), suppose the certificate supplies
an integer shift \(D_q\), a threshold \(A\), and positive constants
\(\eta_{\mathrm{mode}}\) and \(\eta_{\mathrm{reserve}}\), with
\[
  a\equiv q\pmod 3,\qquad m=\frac{2a+D_q}{3}.
\]

\begin{lemma}[Adjacent-to-global hub-off margin]
Let \(F(x)=\sum_k F_kx^k\) have nonnegative unimodal coefficients and mode
at \(m\).  If
\[
  F_m-F_{m-1}\ge \eta F_m,\qquad
  F_m-F_{m+1}\ge \eta F_m,
\]
then for every \(k\ne m\),
\[
  F_m-F_k\ge \eta F_m.
\]
\end{lemma}

\begin{lemma}[Fugacity perturbation]
Let
\[
  \lambda_0=\frac{F_{m-1}}{F_m},\qquad
  \lambda=\frac{F_{m-1}+G_{m-1}}{F_m+G_m}.
\]
Assume \(3/4\le\lambda_0\le 2\), \(0\le T\le 1/2\), and
\[
  0\le G_{m-1}\le TF_m,\qquad 0\le G_m\le TF_m.
\]
Then
\[
  |\lambda-\lambda_0|\le 4T,\qquad \lambda\ge 1/2.
\]
\end{lemma}

\begin{lemma}[Mean perturbation]
Let \(H(x)=\sum_{k=0}^N h_kx^k\) have nonnegative coefficients and
\(H(\lambda)>0\) on the interval between \(\lambda_0\) and \(\lambda\).
If \(\lambda_0,\lambda\ge 1/2\), then
\[
  |\mu_H(\lambda)-\mu_H(\lambda_0)|
  \le \frac{N^2}{2}|\lambda-\lambda_0|.
\]
\end{lemma}

\begin{theorem}[Fixed-\(r\) certificate criterion]
Assume:
\begin{enumerate}
  \item exact finite verification of the Route-2 inequality for \(1\le a<A\);
  \item the adjacent \(F\)-mode margins above for all \(a\ge A\);
  \item hub-on domination strong enough that \(G\) cannot move the certified
        mode \(m\);
  \item hub-off reserve
        \[
          \mu_{F^-_{a,r}}(\lambda_0)
          \ge m-\frac43+\frac{\eta_{\mathrm{reserve}}}{a};
        \]
  \item hub-on perturbation bounded by
        \[
          |\mu_{B_{a,r}}(\lambda)-\mu_{F^-_{a,r}}(\lambda_0)|
          \le \frac{\eta_{\mathrm{reserve}}}{a}.
        \]
\end{enumerate}
Then for every \(a\ge1\),
\[
  \mu_{B_{a,r}}(\lambda)\ge m-\frac32.
\]
\end{theorem}

\section{Expected Output}

The first-pass Lean file in this packet formalizes the adjacent-margin and
fugacity-perturbation arithmetic.  The next proof-agent task is to formalize the
finite-support Gibbs mean-shift argument and then state the full certificate
criterion in Lean.

\end{document}
